\paragraph{Stylized facts and parametric volatility.}
The canonical set of stylized facts \cite{cont2001empirical} — heavy tails \cite{mandelbrot1963variation}, absence of linear autocorrelation \cite{fama1970efficient}, and slow decay of the absolute-return ACF \cite{schwert1989does} — must be reproduced jointly for a generator to be useful downstream.
ARCH/GARCH \cite{engle1982autoregressive, bollerslev1986generalized} encodes volatility persistence through a single decay parameter but imposes a unimodal innovation distribution and, as we confirm, fails distributional tests on equity returns.
Jump-diffusion models \cite{merton1976option, kou2002jump} generate heavy tails but lack a clustering mechanism: after a jump the conditional variance returns immediately to baseline.

\paragraph{Regime-switching and the Ryd\'{e}n et al.\ limitation.}
\citet{hamilton1989new} introduced Markov-switching models to economics, and \citet{schaller1997regime} provided an early demonstration that CRSP monthly equity returns are well-described by a two-state Markov-switching specification with strongly different mean and variance between regimes.
The seminal evaluation of HMMs against the Cont stylized facts is due to \citet{ryden1998stylized}, who fitted Gaussian-emission HMMs with $K = 2$--$3$ states to daily Swedish equity returns and showed that the geometric sojourn-time distribution of a standard Markov chain produces exponential ACF decay too fast to match the observed slow decay.
\citet{bulla2006stylized} restored fidelity with hidden semi-Markov models (HSMMs) that replace the geometric sojourn distribution with flexible alternatives.
\citet{alswaidan2026hybrid} instead kept the discrete HMM framework and introduced a Poisson jump-duration process to force dwell time in tail states.

Within the same family, \citet{nguyen2018hidden} applies a four-state continuous HMM to monthly S\&P~500 prices for out-of-sample trading signal generation, selecting the state count through the AIC, BIC, HQC, and CAIC information criteria.
We adopt the same four-criterion selection (Section~\ref{sec:model_selection}), in addition to our multi-metric distributional score, to place the $K$-sweep recommendation on standard statistical footing.

The present work resolves the Ryd\'{e}n limitation through a third route: increase the state resolution, keep the standard HMM framework, and let the Baum-Welch algorithm learn heterogeneous self-transition probabilities across regimes.
At $K \geq 3$ with quantile-based initialization, high-volatility states develop enough self-persistence that the resulting mixture of $K$ geometric sojourn-time distributions approximates slow ACF decay, effectively recovering the HSMM behavior within an ordinary Markov chain.

\paragraph{Discrete HMM with jumps.}
The immediate predecessor of this paper, \citet{alswaidan2026hybrid}, discretizes returns into Laplace quantile bins, estimates $\mathbf{T}$ by frequency counting, fits bin-conditional Student-$t$ emissions, and adds a Poisson jump-duration process with two grid-calibrated hyperparameters ($\epsilon, \lambda$).
The continuous model presented here eliminates all three sources of complexity: no binning, no separate emission fit (Baum-Welch jointly learns transitions and emissions), and no jump mechanism at moderate $K$.

\paragraph{Cross-asset dependence.}
For multi-asset synthesis, the Single Index Model \cite{sharpe1963simplified} propagates a market factor through a rank-one linear decomposition but distorts each non-market marginal; Sklar's theorem \cite{sklar1959fonctions, nelsen2006introduction} supports a cleaner decomposition in which each asset's fitted marginal is kept and only the copula is estimated.
Gaussian and Student-t copulas \cite{demarta2005tcopula, mcneil2015quantitative} are the standard choices, and rank-reordering simulation \cite{iman1982distribution} injects the copula dependence onto pre-simulated marginal paths without distorting any marginal.
We confirm, on continuous CHMM marginals, the copula-ahead-of-SIM ordering reported by \citet{alswaidan2026hybrid} for discrete HMM marginals, which suggests that the ordering is intrinsic to the rank-reordering construction rather than specific to the marginal model.

\paragraph{Deep generative models and quality assessment.}
GAN-based generators reproduce the visual appearance of return series but typically fail the clustering test unless architectures encode temporal dependence explicitly \cite{takahashi2019modeling, kwon2024can}; neural baselines in \citet{alswaidan2026hybrid} close most of the ACF gap but suffer variance collapse, indicating that the distribution--clustering tension remains in tension for neural approaches.
For evaluation, we adopt the category framework of \citet{stenger2024thinking} and the financial-synthesis emphasis of \citet{assefa2020generating}: no single metric is decisive, so we report a seven-metric panel spanning distributional, tail, and temporal fidelity.
